\documentclass[10pt,a4paper]{article}
\usepackage{amsmath,amssymb,amsthm,graphicx}
\usepackage{amsmath,amssymb,amsthm,xcolor,bm}
\usepackage{hyperref}
\usepackage[nameinlink,noabbrev,capitalize]{cleveref}
\allowdisplaybreaks
\newcommand{\eps}{\varepsilon}
\newcommand{\cP}{\mathcal P}
\newcommand{\AC}{\texttt{AC}}
\newcommand{\poly}{\texttt{poly}}
\newtheorem{theorem}{Theorem}
\newtheorem{conjecture}[theorem]{Conjecture}
\newtheorem{example}[theorem]{Example}
\newtheorem{remark}[theorem]{Remark}
\newtheorem{rem}[theorem]{Remark}
\newtheorem{fact}[theorem]{Fact}
\newtheorem{lemma}[theorem]{Lemma}
\newtheorem{problem}[theorem]{Problem}
\newtheorem{proposition}[theorem]{Proposition}
\newtheorem{corollary}[theorem]{Corollary}
\newtheorem{observation}[theorem]{Observation}
\newtheorem{THM}[theorem]{Theorem}
\theoremstyle{definition}
\newtheorem{definition}[theorem]{Definition}
\crefname{THM}{Theorem}{Theorems}
\Crefname{THM}{Theorem}{Theorems}
\crefname{equation}{}{}
\crefrangeformat{equation}{(#3#1#4) --~(#5#2#6)}
\crefmultiformat{equation}{(#2#1#3)}{, (#2#1#3)}{, (#2#1#3)}{, (#2#1#3)}
\title{Fourier-level lower bound for every low-degree uniform approximator on fractional subcubes}
\author{Eshan Chattopadhyay, Jason Gaitonde, Chin Ho Lee, Shachar Lovett, Abhishek Shetty}
\date{doi:10.4230/LIPIcs.CCC.2021.10; CC BY 4.0}
\begin{document}
\maketitle
\noindent\textit{Editorial note.} Editable transcription of the official CC BY author approximation-converse result and full bearing Sections2.1/3/3.1 context. Independent public copy blanks all unselected material while retaining native original locator lines; formal chapter, statement and equation offsets restored. Author mathematical statements and proofs unchanged. Standard Chebyshev extremal theorem cited rather than recursively reproved. All production and validation solely on Yupeng.\par\medskip
\noindent\textit{Allowed standard prerequisites (editorial).}\begin{itemize}
\item Boolean Fourier analysis, univariate Taylor theorem, standard bounded-wise independence/small-bias generators and prior polarizing-walk amplification. Complete originalSection2 definitions and cited standard amplificationTheorem14 included; all source-local Fourier approximation and random-restriction proofs retained
\end{itemize}\medskip
\setcounter{section}{1}\setcounter{theorem}{6}\setcounter{equation}{1}\section{Preliminaries}
\label{sec:preliminaries}
As in \cite{CHHL} and \cite{CHLT}, we study PRGs for classes $\mathcal{F}$ of $n$-variate Boolean functions that are closed under restriction (that is, fixing any subset of the input variables of a function in the class yields a function that remains in the class). 


\subsection{Fourier Analysis} \label{subsec:fourier}
We briefly recall basic Fourier analysis: any Boolean function $f:\{-1,1\}^n \to \{-1,1\}$ admits a unique multilinear expansion, also known as the \emph{Fourier expansion}, given by 
\begin{equation}
\label{eq:fourierexpansion}
    f(\mathbf{x})=\sum_{S\subseteq [n]} \hat{f}(S)\mathbf{x}^S,
\end{equation}
where we write $\mathbf{x}^S\triangleq \prod_{i\in S} x_i$. The   Fourier coefficient $\hat{f}(S)$ is given by
\begin{equation*}
    \hat{f}(S)=\mathbb{E}_{\mathbf{X}\sim \{-1,1\}^n}[f(\mathbf{X})\mathbf{X}^S].
\end{equation*}
For more on Fourier analysis of Boolean functions, see the excellent book by O'Donnell \cite{ODonnelBook}. One may thus extend the domain of $f$ to $[-1,1]^n$, where $f(\mathbf{x})$ for arbitrary $\mathbf{x}$ is evaluated according to the expression in \Cref{eq:fourierexpansion}. Note that in this case, $f(\mathbf{0})=\hat{f}(\emptyset)=\mathbb{E}_{\mathbf{U}_n}[f(\mathbf{U}_n)]$. One of the main parameters of interest from the Fourier expansion for this framework is the following:


\begin{definition}
The \emph{level-$k$ mass of a Boolean function} $f$ is 
\begin{equation*}
    L_{1,k}(f)\triangleq \sum_{S\subseteq [n]:\vert S\vert=k} \vert \hat{f}(S)\vert,
\end{equation*}
and the \emph{level-$k$ mass of a class $\mathcal{F}$} is
$L_{1,k}(\mathcal{F})\triangleq \max_{f\in \mathcal{F}} L_{1,k}(f)$.
\end{definition}

In this work, we will show how to construct PRGs whose seed length depends on the following, smaller quantity:

\begin{definition}
\label{def:mk}
For any multilinear polynomial $f:\mathbb{R}^n\to \mathbb{R}$ given by $f(\mathbf{x})=\sum_{S\subseteq [n]} \hat{f}(S)\mathbf{x}^S$, define the level-$k$ part by
\begin{equation*}
    f_{k}(\mathbf{x})\triangleq\sum_{S\subseteq [n]:\vert S\vert=k} \hat{f}(S)\mathbf{x}^S,
\end{equation*}
and further define $f_{<k}(\mathbf{x})\triangleq \sum_{i=0}^{k-1} f_i(\mathbf{x})$ and $f_{\geq k}(\mathbf{x})\triangleq \sum_{i=k}^{n} f_i(\mathbf{x})$. 
Then we define the \emph{level-$k$ absolute Fourier sum of $f$} by
\begin{equation*}
    M_k(f)\triangleq \max_{\mathbf{x}\in [-1,1]^n}\bigg\vert \sum_{S:\vert S\vert=k} \hat{f}(S)\mathbf{x}^S\bigg\vert =  \max_{\mathbf{x}\in \{-1,1\}^n}\bigg\vert \sum_{S:\vert S\vert=k} \hat{f}(S)\mathbf{x}^S\bigg\vert
\end{equation*}
and analogously define $M_k(\mathcal{F})\triangleq \max_{f\in \mathcal{F}} M_k(f)$ for a class $\mathcal{F}$.
\end{definition}
Note that the equality arises by multilinearity, and clearly we have $M_k(f)\leq L_{1,k}(f)$ by the triangle inequality. Without loss of generality, we may further assume that our class is closed under flipping the image, i.e. we may suppose that $f\in \mathcal{F}$ if and only if $-f\in \mathcal{F}$; this transformation does not change either $L_{1,k}(f)$ or $M_k(f)$, and therefore the same bound on the class still holds when completing it to include all such functions. If this is the case, we get the more striking identity:

\begin{lemma}
\label{lem:unsigned}
Suppose that $\mathcal{F}$ is closed under negation of variables and that $f\in \mathcal{F}$ implies $-f\in \mathcal{F}$. Then
\begin{equation*}
    M_k(\mathcal{F})=\max_{f\in \mathcal{F}} \sum_{S:\vert S\vert = k} \hat{f}(S) = \max_{f \in  \mathcal{F} } f_k \left( \mathbf{1} \right) .
\end{equation*}
\end{lemma}

To see why this holds, simply note that if $(f,\mathbf{z})\in \mathcal{F}\times \{-1,1\}^n$ is a maximizer in the definition of $M_k(\mathcal{F})$ (where we may now assume that the sign is positive), then by replacing the function $f(\mathbf{x})$ with $g(\mathbf{x})=f(\mathbf{x}\circ \mathbf{z})$, where $\circ$ denotes componentwise multiplication, we have 
\begin{equation*}
    M_k(\mathcal{F})
    = \Biggl| \sum_{S:\vert S\vert=k} \hat{f}(S)\mathbf{z}^S \Biggr|
    = \sum_{S:\vert S\vert=k} \hat{g}(S)=\max_{h\in \mathcal{F}} \sum_{S:\vert S\vert=k} \hat{h}(S).
\end{equation*}
In particular, it suffices to bound the \emph{unsigned level-$k$ Fourier sum} of such a class.

Lastly, we require the following notion:

\begin{definition}
Let $\mathcal{F}$ be a class of $n$-variate multilinear polynomials that is closed under restrictions.  Define $\mathrm{conv}(\mathcal{F})$ as the convex closure of $\mathcal{F}$,
\begin{equation*}
    \mathrm{conv}(\mathcal{F})\triangleq \left\{\sum_{f\in \mathcal{F}} \lambda_f f\,\bigg\vert \sum_{f\in \mathcal{F}} \lambda_f=1, \lambda_f\geq 0 \,\,\forall f\in \mathcal{F}\right\}.
\end{equation*}
\end{definition}
We briefly note the following two elementary facts: first, by the assumption that $\mathcal{F}$ is closed under restrictions, the same is true of $\mathrm{conv}(\mathcal{F})$. The second is the following simple claim:

\begin{lemma}
For any class $\mathcal{F}$ of Boolean functions, $M_k(\mathcal{F})=M_k(\mathrm{conv}(\mathcal{F}))$.
\end{lemma}
\begin{proof}
One direction is obvious: as $\mathcal{F}\subseteq \mathrm{conv}\mathcal{F}$, clearly $M_k(\mathcal{F})\leq M_k(\mathrm{conv}(\mathcal{F}))$. In the other direction, let $g=\sum_{f\in \mathcal{F}} \lambda_f f$ be an arbitrary element of $\mathrm{conv}(\mathcal{F})$, where $\lambda_f\geq 0$ and $\sum_{f\in \mathcal{F}}\lambda_f=1$. Then
\begin{align*}
    M_k(g) &= \max_{\mathbf{x}\in \{-1,1\}^n} \Biggl| \sum_{S\subseteq [n]:\vert S\vert=k} \widehat{g}(S) \mathbf{x}^S \Biggr| \\
    &=\max_{\mathbf{x}\in \{-1,1\}^n} \Biggl| \sum_{S\subseteq [n]:\vert S\vert=k} \Bigl( \sum_{f\in \mathcal{F}}\lambda_f\widehat{f}(S) \Bigr) \mathbf{x}^S \Biggr| \\
    &\leq \sum_{f\in \mathcal{F}} \lambda_f \max_{\mathbf{x}\in \{-1,1\}^n} \Biggl| \sum_{S\subseteq [n]:\vert S\vert=k} \widehat{f}(S) \mathbf{x}^S \Biggr| \\
    &\leq \max_{f\in \mathcal{F}} M_k(f).
\end{align*}
The reverse inequality immediately follows.
\end{proof}


\setcounter{theorem}{16}\section{Low-Degree Polynomial Approximations on Subcubes}
\label{sec:lowdeg}

Throughout this section, we assume that $\mathcal{F}$ is a class of $n$-variate Boolean functions closed under restrictions.  As mentioned above, the main result from which we derive our improvements in constructing pseudorandom generators is essentially a statement about low-degree polynomial approximations on subcubes $[-c,c]^n$ for $c<1$.   We remark that this setting is equivalent to approximating \emph{noisy} versions $T_cf$ on $[-1,1]^n$, where $T_{\rho}$ is the $\rho$-noise operator.  This is because for any $\mathbf{y} \in [-c,c]^n$, we can write $\mathbf{y} = c\mathbf{x}$ for some $\mathbf{x} \in [-1,1]^n$ and
\[
f(\mathbf{y})
=f(c\mathbf{x})
= \sum_{S \subseteq [n]} \hat{f}(S) (c\mathbf{x})^S
= \sum_{S \subseteq [n]} c^{|S|} \hat{f}(S) \mathbf{x}^S
= T_cf(\mathbf{x}) .
\]
In general, given any $k\leq n$, $c\geq 0$, and any $f\in \mathcal{F}$, let $\eps_{c,k}(f)$ be defined by
\begin{equation}
\label{eq:fapprox}
    \eps_{c,k}(f) \triangleq \inf_{g: \mathrm{deg}(g)< k}\max_{\mathbf{x}\in [-c,c]^n}  \vert f(\mathbf{x})-g(\mathbf{x})\vert,
\end{equation}
and extend the definition to function classes by
\begin{equation*}
    \eps_{c,k}(\mathcal{F}) \triangleq \max_{f\in \mathcal{F}} \eps_{c,k}(f).
\end{equation*}
Now, given $\eps>0$, $k\leq n$, and the class $\mathcal{F}$, define $c_k(\mathcal{F},\eps)$ by
\begin{equation*}
    c_k(\eps,\mathcal{F})\triangleq \max\{ c\geq 0: \eps_{c,k}(\mathcal{F})\leq \eps\}.
\end{equation*}
In words, $c_k(\eps,\mathcal{F})$ measures how small a hypercube we must take to ensure that for every function in our class, there exists a degree-($k-1$) approximating polynomial that agrees with $f$ up to a uniform $\eps$ error on the subcube $[-c,c]^n$; by multilinearity, it actually suffices that this holds at the extreme points $\{-c,c\}^n$. Note that \cref{eq:fapprox} can be formulated as a linear program and its optimal solution is the best low-degree  $\ell_{\infty}$-approximation to $f$.


The main technical claim in this section is that we bound $c_k \left( \eps , \mathcal{F} \right)$ in terms of $M_k(\mathcal{F})$.  Specifically, we show that for any class $\mathcal{F}$ that is closed under restrictions, truncating the Fourier expansion of a function $f \in \mathcal{F}$ to its first $(k-1)$ levels serves as a good approximation to $f$ on a sufficiently small hypercube around the origin.

\begin{THM}
\label{thm:main}
        Let $f\in \mathcal{F}$ that is closed under restrictions. Then for all $c \in (0,1)$, we have     \begin{equation*}
        \max_{\mathbf{x} \in [-c , c]^n } \left \vert f_{\geq k} (\mathbf{x}) \right\vert \leq \left( \frac{c}{1 - c} \right)^k M_k(\mathcal{F}).
    \end{equation*}
    In particular, it follows that
    \begin{equation*}
        \eps_{c,k}(\mathcal{F})\leq \left( \frac{c}{1 - c} \right)^k M_k(\mathcal{F}).
    \end{equation*}
\end{THM}
From \Cref{thm:main}, one immediately obtains a lower bound on $c_k(\eps,\mathcal{F})$:
\begin{corollary}
\label{cor:ckbound}
For any class $\mathcal{F}$ that is closed under restrictions, and any $\eps>0$ and $k\leq n$, 
\begin{equation*}
    c_k(\eps,\mathcal{F}) = \Omega\left(\left(\frac{\eps}{M_k(\mathcal{F})}\right)^{1/k}\right)
\end{equation*}
\end{corollary}
\begin{proof}
Observe that by setting $c=\Omega\left(\left(\frac{\eps}{M_k(\mathcal{F})}\right)^{1/k}\right)$ in \Cref{thm:main}, the right side is bounded by $\eps$. Because $f_{\geq k}=f-f_{<k}$ and $f_{<k}$ has degree strictly less than $k$, it follows immediately from the definition of $c_k(\eps,\mathcal{F})$ that $c_k(\eps,\mathcal{F})$ is at least $c$.
\end{proof}
    We now return to the proof of \Cref{thm:main}. To prove this result, we require the following intermediate claims. The first simply shows that we may always bound the contribution of the level-$k$ part of any function in $\mathcal{F}$ by simply rescaling the argument: 
    \begin{lemma} 
    \label{lem:mult}
        Let $f \in \mathrm{conv}(\mathcal{F})$. 
        Then, for all $c\in ( 0,1 )$ and $\mathbf{x} \in [-c , c  ]^n$, we have 
        \begin{equation*}
            \vert f_k(\mathbf{x})\vert \leq c^k M_k(\mathcal{F}). 
        \end{equation*}
    \end{lemma}
    
    \begin{proof}
    Observe that $c^{-1}\mathbf{x}\in [-1,1]^n$ by assumption, and by homogeneity of $f_{k}$ as a polynomial, we have
    \begin{equation*}
        \vert f_k(\mathbf{x})\vert=c^k \vert f_k(c^{-1}\mathbf{x})\vert\leq c^k M_k(\mathrm{conv}(\mathcal{F}))= c^k M_k(\mathcal{F}). \qedhere
    \end{equation*}
    \end{proof}

The next simple yet powerful claim shows that one can ``recenter'' functions in $\mathcal{F}$ and they remain in $\mathrm{conv}(\mathcal{F})$ (and therefore, enjoy the same Fourier bounds). This random restriction technique is a key tool in \cite{CHHL}.

    \begin{lemma} \label{lem:conv}
        Let $f \in \mathrm{conv}(\mathcal{F})$, $\mathbf{a} \in [-1,1]^n$ and $\mathbf{b} \in [0,1]$ such that $ \vert a_i\vert + b_i \leq 1$ for all $i\in [n]$. 
        Define $\tilde{f} $ by $\tilde{f}(\mathbf{x}) = f(\mathbf{a} + \mathbf{b} \circ \mathbf{x})$, where $\circ$ denotes componentwise multiplication. Then, $\tilde{f}  \in \mathrm{conv}(\mathcal{F})$. 
    \end{lemma}
    \begin{proof}
        Given $\mathbf{a}$ and $\mathbf{b}$, define a distribution $D_i$ on $Z_i= \{ -1,1 , x_i\} $ where $x_i$ is treated as formal variable, such that $\mathbb{E}_{y_i \sim D_i }[ y_i] = a_i + b_i x_i  $; note that this is possibly by the assumption that $\vert a_i\vert + b_i \leq 1$. 
        Let $D = \prod_i D_i  $ be the product distribution of the $D_i$. For any $\mathbf{z} \in \prod_i Z_i $, define $f_{\mathbf{z}}(\mathbf{x})$ as the function obtained by setting $x_i=z_i$ for each $i$; in particular, each variable gets set to $\pm 1$ or remains a formal variable.  By our assumption on the closure of $\mathcal{F}$, we clearly have $f_{\mathbf{z}}\in \mathcal{F}$ for any $\mathbf{z}$. By multilinearity and independence of the product distribution, we have $f(\mathbf{a} + \mathbf{b} \circ \mathbf{x}) = \mathbb{E}_{\mathbf{z} \sim D}[ f_{\mathbf{z}}(\mathbf{x})]  $. Thus $\tilde{f} \in \mathrm{conv}(\mathcal{F})$. 
    \end{proof}

As mentioned before, our approach will be to bound the higher-order terms of the Fourier expansion at the fractional points of the fractional PRG via the error term that arises in Taylor's theorem. Denote by $h^{(k)}$ the $k$-th derivative of any $C^k$ function $h:\mathbb{R}\to \mathbb{R}$. We then have the following claim:
    \begin{lemma} 
    \label{lem:derivative}
        Let $f:\mathbb{R}^n \to \mathbb{R}  $ be multilinear and let $\mathbf{x} \in \mathbb{R}^n $. 
        Define $g:\mathbb{R} \to \mathbb{R}$ by $g(t) = f(t \mathbf{x}) $. 
        Then, 
        \begin{equation*}
            g^{(k)}(0) = k! \cdot f_k ( \mathbf{x}). 
        \end{equation*}
    \end{lemma}
    \begin{proof}
        From the definition, it follows that 
        \begin{equation*}
            g( t) = \sum_{S \subseteq [ n ] } t^{\vert S\vert} \hat{f}(S) \mathbf{x}^S.
         \end{equation*}
         Differentiating $g$ with respect to $t$, we get 
         \begin{equation*}
            g^{(k)}(t)=\sum_{S: \vert S\vert  \geq k} \Bigl(\, \prod_{i=0}^{k-1}( \vert S\vert-i) \Bigr) t^{ \vert S\vert -k} \hat{f}(S) \mathbf{x}^{S}. 
         \end{equation*}
         Setting $t = 0$ eliminates all of the monomials with $\vert S\vert > k $, giving us the required bound. 
    \end{proof}

The last intermediate result we require connects the function defined in the previous part with our assumed Fourier bounds:
    \begin{lemma}
    \label{lemma:finallem}
        Let $f \in \mathrm{conv}(\mathcal{F}) $, $c \in (0,1) $ and $\mathbf{x} \in [ -c , c]^n $. 
        Define $g$ as in \cref{lem:derivative}. 
        Then, 
        \begin{equation*}
            \max _{s \in[0,1]} \bigl| g^{(k)}(s) \bigr|
            \leq\left(\frac{c}{1-c}\right)^{k} \cdot k! \cdot M_{k}(\mathcal{F})
        \end{equation*}
    \end{lemma}
    \begin{proof}
        Fix $s \in [ 0,1]$ and let $\lambda=1-c \in [0,1]$.  Define the auxiliary function $\tilde{f}(\mathbf{y})=f(s\mathbf{x}+\lambda\mathbf{y})$. Writing $\mathbf{a}=s\mathbf{x}$ and $\mathbf{b}=(\lambda,\ldots,\lambda)$, we clearly have $s |x_i| + \lambda \leq 1$, so we may apply \cref{lem:conv} to see that $\tilde{f}\in \mathrm{conv}(\mathcal{F})$.
        Now writing $\tilde{g}(t) = \tilde{f}(t\mathbf{x})=f(s\mathbf{x}+\lambda t\mathbf{x})$, we also have $\tilde{g}(t)=g(s+t\lambda)$. By the chain rule, differentiating both sides $k$ times and then setting $t=0$, we have
        \begin{equation*}
            \lambda^k g^{(k)}(s)=\tilde{g}^{(k)}(0).
        \end{equation*}
        On the other hand, by \cref{lem:derivative}, we have $\tilde{g}^{(k)}(0)=k!\cdot \tilde{f}_k(\mathbf{x})$, and as $\tilde{f}\in \mathrm{conv}(\mathcal{F})$ by \cref{lem:conv}, we conclude using \cref{lem:mult} that
        \[
            \bigl|g^{(k)}(s)\bigr|=\left|\frac{\tilde{g}^{(k)}(0)}{\lambda^k}\right|\leq \left(\frac{c}{1-c}\right)^k \cdot k!\cdot M_k(\mathcal{F}) . \qedhere
        \]
    \end{proof}
    With these intermediate claims taken care of, we may now put them together to obtain \Cref{thm:main}.
    \begin{proof}[Proof of \cref{thm:main}]
    The second statement follows immediately from the first by setting $g=f_{<k}$ for any given $f$, and noticing that $f-g=f_{\geq k}$. Therefore, we focus on the first statement.
    
Let $f\in \mathcal{F}, \mathbf{x}\in [-c,c]^n$ and define $g(t)=f(t\mathbf{x})$. Then, by Taylor expanding $g$ about $t=0$ and evaluating $g$ at $t=1$, we have

\begin{align} \label{eq:decomp-g}
    g(1)=\sum_{i<k} \frac{g^{(i)}(0)}{i!} +R_k, 
\end{align}
where $R_k$ is the error term and is given in Lagrange form by
\begin{equation*}
    R_k = \frac{g^{(k)}(s)}{k!}
\end{equation*}
for some $s\in (0,1)$. By \cref{lem:derivative}, we easily see that the first term in the right hand side of \cref{eq:decomp-g} is precisely $f_{<k}(\mathbf{x})$, and as $g(1)=f(\mathbf{x})$, we clearly then must have $R_k = f_{\geq k}(\mathbf{x})$.  Therefore, by \cref{lemma:finallem}, we obtain
\begin{equation*}
    \vert f_{\geq k}(\mathbf{x})\vert = \left\vert\frac{g^{(k)}(s)}{k!}\right\vert \leq \left(\frac{c}{1-c}\right)^k M_k(\mathcal{F}),
\end{equation*}
as desired.
\end{proof}

\subsection{Lower Bounds via Chebyshev Polynomials}

In this subsection, we show that our bounds on the uniform error of any low-degree polynomial approximator are essentially tight for a reasonable range of $c<1$. Recall that \Cref{thm:main} shows that the low-degree Fourier expansion is an excellent approximator to the original function for $c$ small enough; we now show that this bound cannot be significantly improved for a reasonable range of $c$ using \emph{any} approximator. Our main result of this section is the following converse:


\begin{THM}
\label{thm:cheby}
Let $\mathcal{F}$ be any class of $n$-variate multilinear functions that are closed under restrictions.  Then for any $c\leq \min\left(\frac{1}{3}, 3^{-k} \frac{M_k(\mathcal{F})}{M_{k+1}(\mathcal{F})}\right)$, we have
\begin{equation*}
    \eps_{c,k}(\mathcal{F}) \geq \left(\frac{c}{2}\right)^k M_k(\mathcal{F}).
\end{equation*}
\end{THM}


Recall that on the interval $[-1,1]$, the Chebyshev polynomials give the minimum $\ell_{\infty}$ norm among all polynomials with same leading coefficient in magnitude:

\begin{fact}[Theorem 1.5.4 of \cite{rahman2002analytic}]
\label{fact:cheb3}
If a polynomial $f:\mathbb{R}\to \mathbb{R}$ is monic of degree $n$, then $\max_{x\in [-1,1]} \vert f(x)\vert\geq 2^{-n+1}$, with equality if and only if $f=T_n$, the normalized $n$-th Chebyshev polynomial.
\end{fact}

\begin{proof}[Proof of \Cref{thm:cheby}]
Let $(f,\mathbf{x})$ attain the maximum in the definition of $M_k(\mathcal{F})$, namely
\begin{equation*}
    M_k(\mathcal{F}) = \left\vert \sum_{S\subseteq [n]:\vert S\vert=k} \widehat{f}(S)\mathbf{x}^S\right\vert.
\end{equation*}
First, note that the claim is trivial if every function in $\mathcal{F}$ is of degree at most $k$, because then $f_{\geq k}$ is a homogeneous polynomial of degree $k$ and this lower bound is trivial. Under this assumption, $M_{k+1}(\mathcal{F})>0$. Fix $c\in (0,1)$ and let $p:[-1,1]^n\to \mathbf{R}$ be any multilinear polynomial of degree strictly less than $k$. Define the univariate function $g:[-1,1]\to \mathbb{R}$ by
\begin{equation*}
    g(t) = f(tc\mathbf{x})- p(tc\mathbf{x}).
\end{equation*}
By taking the Fourier expansion of $f$, it is easy to see that the coefficient of $t^{\ell}$ for $\ell\geq k$ is precisely
\begin{equation*}
    c^{\ell}\sum_{S\subseteq [n]:\vert S\vert=\ell} \widehat{f}(S)\mathbf{x}^S,
\end{equation*}
so that the coefficient of $t^k$ is equal to $c^k M_k(\mathcal{F})$ in magnitude.
We then have
\begin{align*}
    \sup_{\mathbf{z}\in [-c,c]^n} \vert f(\mathbf{z})-p(\mathbf{z})\vert&\geq \max_{\mathbf{z}\in [-c\mathbf{x},c\mathbf{x}]} \vert f(\mathbf{z})-p(\mathbf{z})\vert\\
    &=\sup_{t\in [-1,1]} \vert g(t)\vert\\
    &\geq \sup_{t\in [-1,1]} \vert g_{\leq k}(t)\vert - \sup_{t\in [-1,1]} \vert g_{\geq k+1}(t)\vert.
\end{align*}
By \Cref{fact:cheb3}, the first term is at least $c^kM_k(\mathcal{F})/2^{k-1}$. On the other hand, the second term can be bounded using \Cref{thm:main} by
\begin{equation*}
    \sup_{t\in [-1,1]} \vert g_{\geq k+1}(t)\vert\leq \left(\frac{c}{1-c}\right)^{k+1}M_{k+1}(\mathcal{F}).
\end{equation*}
Therefore, we obtain
\begin{equation*}
    \sup_{\mathbf{z}\in [-c,c]^n} \vert f(\mathbf{z})-p(\mathbf{z})\vert\geq 2\left(\frac{c}{2}\right)^kM_k(\mathcal{F})-\left(\frac{c}{1-c}\right)^{k+1}M_{k+1}(\mathcal{F}).
\end{equation*}
It is straightforward to verify that for $c\leq \min\left(1/3, 3^{-k}  \frac{M_k(\mathcal{F})}{M_{k+1}(\mathcal{F})}\right)$, the second term is bounded by half of the first. Because $p$ was an arbitrary low-degree multilinear polynomial, the claim follows.
\end{proof}

\newcommand{\etalchar}[1]{$^{#1}$}
\begin{thebibliography}{DFOC{\etalchar{+}}11}

\bibitem[ACK{\etalchar{+}}20]{DBLP:conf/stacs/Arunachalam0KSW20}
Srinivasan Arunachalam, Sourav Chakraborty, Michal Kouck{\'{y}}, Nitin Saurabh,
  and Ronald de~Wolf.
\newblock Improved bounds on {F}ourier entropy and min-entropy.
\newblock In Christophe Paul and Markus Bl{\"{a}}ser, editors, {\em 37th
  International Symposium on Theoretical Aspects of Computer Science, {STACS}
  2020, March 10-13, 2020, Montpellier, France}, volume 154 of {\em LIPIcs},
  pages 45:1--45:19. Schloss Dagstuhl - Leibniz-Zentrum f{\"{u}}r Informatik,
  2020.

\bibitem[BS20]{DBLP:journals/corr/abs-2008-07003}
Nikhil Bansal and Makrand Sinha.
\newblock {$k$}-{F}orrelation {O}ptimally {S}eparates {Q}uantum and {C}lassical
  {Q}uery {C}omplexity.
\newblock {\em CoRR}, abs/2008.07003, 2020.

\bibitem[CHH{\etalchar{+}}20]{xorlemma}
Eshan Chattopadhyay, Pooya Hatami, Kaave Hosseini, Shachar Lovett, and David
  Zuckerman.
\newblock {XOR} lemmas for resilient functions against polynomials.
\newblock In {\em Proceedings of the 52nd Annual ACM SIGACT Symposium on Theory
  of Computing}, STOC 2020, page 234–246, New York, NY, USA, 2020.
  Association for Computing Machinery.

\bibitem[CHHL19]{CHHL}
Eshan Chattopadhyay, Pooya Hatami, Kaave Hosseini, and Shachar Lovett.
\newblock Pseudorandom generators from polarizing random walks.
\newblock {\em Theory of Computing}, 15(10):1--26, 2019.

\bibitem[CHLT19]{CHLT}
Eshan Chattopadhyay, Pooya Hatami, Shachar Lovett, and Avishay Tal.
\newblock {Pseudorandom Generators from the Second {F}ourier Level and
  Applications to {AC0} with Parity Gates}.
\newblock In Avrim Blum, editor, {\em 10th Innovations in Theoretical Computer
  Science Conference (ITCS 2019)}, volume 124 of {\em Leibniz International
  Proceedings in Informatics (LIPIcs)}, pages 22:1--22:15, Dagstuhl, Germany,
  2019. Schloss Dagstuhl--Leibniz-Zentrum fuer Informatik.

\bibitem[CHRT18]{chattopadhyay2018improved}
Eshan Chattopadhyay, Pooya Hatami, Omer Reingold, and Avishay Tal.
\newblock Improved pseudorandomness for unordered branching programs through
  local monotonicity.
\newblock In {\em Proceedings of the 50th Annual {ACM} {SIGACT} Symposium on
  Theory of Computing, {STOC} 2018, Los Angeles, CA, USA, June 25-29, 2018},
  pages 363--375, 2018.

\bibitem[DFOC{\etalchar{+}}11]{defant2011bohnenblust}
Andreas Defant, Leonhard Frerick, Joaquim Ortega-Cerd{\`a}, Myriam
  Ouna{\"\i}es, and Kristian Seip.
\newblock The {B}ohnenblust-{H}ille inequality for homogeneous polynomials is
  hypercontractive.
\newblock {\em Annals of mathematics}, pages 485--497, 2011.

\bibitem[GRZ20]{RazXOR}
Uma Girish, Ran Raz, and Wei Zhan.
\newblock Lower bounds for {XOR} of forrelations.
\newblock {\em CoRR}, abs/2007.03631, 2020.

\bibitem[GSTW16]{GSTW}
Parikshit Gopalan, Rocco~A. Servedio, Avishay Tal, and Avi Wigderson.
\newblock Degree and sensitivity: tails of two distributions, 2016.

\bibitem[GSW16]{GSW16}
Parikshit Gopalan, Rocco~A. Servedio, and Avi Wigderson.
\newblock Degree and sensitivity: Tails of two distributions.
\newblock In Ran Raz, editor, {\em 31st Conference on Computational Complexity,
  {CCC} 2016, May 29 to June 1, 2016, Tokyo, Japan}, volume~50 of {\em LIPIcs},
  pages 13:1--13:23. Schloss Dagstuhl - Leibniz-Zentrum f{\"{u}}r Informatik,
  2016.

\bibitem[Lee19]{DBLP:conf/coco/Lee19}
Chin~Ho Lee.
\newblock Fourier bounds and pseudorandom generators for product tests.
\newblock In Amir Shpilka, editor, {\em 34th Computational Complexity
  Conference, {CCC} 2019, July 18-20, 2019, New Brunswick, NJ, {USA}}, volume
  137 of {\em LIPIcs}, pages 7:1--7:25. Schloss Dagstuhl - Leibniz-Zentrum
  f{\"{u}}r Informatik, 2019.

\bibitem[LMN89]{LMN89}
Nathan Linial, Yishay Mansour, and Noam Nisan.
\newblock Constant depth circuits, {F}ourier transform, and learnability.
\newblock In {\em 30th Annual Symposium on Foundations of Computer Science},
  pages 574--579. IEEE, 1989.

\bibitem[Mon12]{montanaro2012some}
Ashley Montanaro.
\newblock Some applications of hypercontractive inequalities in quantum
  information theory.
\newblock {\em Journal of Mathematical Physics}, 53(12):122206, 2012.

\bibitem[NN90]{naor1993small}
Joseph Naor and Moni Naor.
\newblock Small-bias probability spaces: Efficient constructions and
  applications.
\newblock In Harriet Ortiz, editor, {\em Proceedings of the 22nd Annual {ACM}
  Symposium on Theory of Computing, May 13-17, 1990, Baltimore, Maryland,
  {USA}}, pages 213--223. {ACM}, 1990.

\bibitem[O'D14]{ODonnelBook}
Ryan O'Donnell.
\newblock {\em Analysis of {B}oolean Functions}.
\newblock Cambridge University Press, 2014.

\bibitem[Raz87]{Razborov1987}
Alexander~A. Razborov.
\newblock Lower bounds on the size of bounded depth circuits over a complete
  basis with logical addition.
\newblock {\em Mathematical notes of the Academy of Sciences of the USSR},
  41(4):333--338, Apr 1987.

\bibitem[RS02]{rahman2002analytic}
Qazi~Ibadu Rahman and Gerhard Schmeisser.
\newblock {\em Analytic theory of polynomials}, volume~26 of {\em London
  Mathematical Society Monographs. New Series}.
\newblock The Clarendon Press, Oxford University Press, Oxford, 2002.

\bibitem[RT19]{RazTal}
Ran Raz and Avishay Tal.
\newblock Oracle separation of {BQP} and {PH}.
\newblock In Moses Charikar and Edith Cohen, editors, {\em Proceedings of the
  51st Annual {ACM} {SIGACT} Symposium on Theory of Computing, {STOC} 2019,
  Phoenix, AZ, USA, June 23-26, 2019}, pages 13--23. {ACM}, 2019.

\bibitem[Smo87]{Smolensky1}
Roman Smolensky.
\newblock Algebraic methods in the theory of lower bounds for {B}oolean circuit
  complexity.
\newblock In {\em Proceedings of the Nineteenth Annual ACM Symposium on Theory
  of Computing}, STOC ’87, page 77–82, New York, NY, USA, 1987. Association
  for Computing Machinery.

\bibitem[Smo93]{Smolensky2}
Roman Smolensky.
\newblock On representations by low-degree polynomials.
\newblock In {\em Proceedings of the 1993 IEEE 34th Annual Foundations of
  Computer Science}, SFCS ’93, page 130–138, USA, 1993. IEEE Computer
  Society.

\bibitem[SSW20]{DBLP:journals/eccc/SherstovSW20}
Alexander~A. Sherstov, Andrey~A. Storozhenko, and Pei Wu.
\newblock An optimal separation of randomized and quantum query complexity.
\newblock {\em Electron. Colloquium Comput. Complex.}, 27:128, 2020.

\bibitem[Tal17]{tal2017tight}
Avishay Tal.
\newblock Tight bounds on the {F}ourier spectrum of {AC0}.
\newblock In Ryan O'Donnell, editor, {\em 32nd Computational Complexity
  Conference, {CCC} 2017, July 6-9, 2017, Riga, Latvia}, volume~79 of {\em
  LIPIcs}, pages 15:1--15:31. Schloss Dagstuhl - Leibniz-Zentrum f{\"{u}}r
  Informatik, 2017.

\bibitem[Tal19]{DBLP:journals/eccc/Tal19}
Avishay Tal.
\newblock Towards optimal separations between quantum and randomized query
  complexities.
\newblock {\em Electron. Colloquium Comput. Complex.}, 26:179, 2019.

\bibitem[Vad12]{vadhan2012pseudorandomness}
Salil~P. Vadhan.
\newblock Pseudorandomness.
\newblock {\em Foundations and Trends in Theoretical Computer Science},
  7(1--3):1--336, 2012.

\bibitem[Vio09]{viola2009sum}
Emanuele Viola.
\newblock The sum of $d$ small-bias generators fools polynomials of degree $d$.
\newblock {\em Computational Complexity}, 18(2):209--217, 2009.

\bibitem[Wu20]{wu2020stochastic}
Xinyu Wu.
\newblock A stochastic calculus approach to the oracle separation of {BQP} and
  {PH}.
\newblock {\em CoRR}, abs/2007.02431, 2020.

\end{thebibliography}
\end{document}
